\section{Baseline Profiling}
\label{sec:profiling}

The telemetry of the first AWS run (Table~\ref{tab:awsbaseline}) splits each file's time into its parts. On average a
50-string file took 21.42~s, of which \textbf{21.27~s (99.3\%) was model inference}. Database lookups, S3 transfers and
XML processing together took about 0.15~s. Any meaningful speed-up therefore has to come from the model stage: either
making each model call cheaper, or making fewer of them.

\subsection{How Batching Worked}
Phase~2 processed one file and one language at a time. It filled a batch with untranslated strings in file order until
the batch held 52 strings or an \emph{estimated} 1,024 tokens (characters divided by four), sent it to the model with
beam size~5 (five candidate translations per string), and waited for the result before sending the next batch.

\subsection{Flaws Found}
\begin{table}[H]
\centering
\caption[Flaws of the original batching]{Flaws of the original batching logic, and the evidence for each.}
\label{tab:flaws}
\begin{tabularx}{\linewidth}{@{}l L@{}}
\toprule
\textbf{Flaw} & \textbf{Evidence} \\
\midrule
Guessed token counts & batches were cut by characters $\div$ 4, not by the tokens the model actually reads \\
Mixed-length batches & short and long strings shared batches; \textbf{30\%} of the model's input positions were padding \\
One batch at a time & only one batch was in the engine at any moment; the two cores were \textbf{59\%} busy on average \\
Five candidates per string & beam~5 decodes five hypotheses for every string, the dominant cost (Chapter~\ref{sec:batching}) \\
No output limit & a single string with an unusually long output could hold up its whole batch \\
\bottomrule
\end{tabularx}
\end{table}

\keyresult{\textbf{RQ1.} Model inference accounts for 99.3\% of the translation stage. Within it, the original batching
wasted input positions on padding (30\%), left part of the CPU idle (59\% busy) and paid for five candidate translations
per string.}
